% Copyright 2026 Adib Hasan. LPPL 1.3c or later.
\documentclass[language=assamese,11pt,oneside,openany]{bookbangla}
% The guide demonstrates the real class, with an A4 layout for documentation.
\setstocksize{297mm}{210mm}
\settrimmedsize{297mm}{210mm}{*}
\settrims{0pt}{0pt}
\setlrmarginsandblock{22mm}{22mm}{*}
\setulmarginsandblock{22mm}{22mm}{*}
\checkandfixthelayout
\linespread{1}
\setmainfont{TeX Gyre Pagella}
\setsansfont{TeX Gyre Heros}
\setmonofont{TeX Gyre Cursor}
\setbanglamainfont{Kalpurush}
% LaTeXbangla registers Assamese; Polyglossia supplies Bengali directly.
% The combined guide uses both languages with the same configured font.
\setotherlanguage[numerals=bengali,changecounternumbering=true]{bengali}
\NewCommandCopy\bengalifont\assamesefont
\newcommand*\bengalifontsf{\bengalifont}
\newcommand*\bengalifonttt{\bengalifont}
\makeatletter
\bookbangla@loadmodule{labels-bangla}
\makeatother
\usepackage{minted,hologo,needspace}
\tcbuselibrary{raster}
\definecolor{BanglaGreen}{HTML}{006847}
\definecolor{BanglaInk}{HTML}{263640}
\definecolor{BanglaRed}{HTML}{C92935}
\definecolor{CodePaper}{HTML}{F4F6F8}
\definecolor{ResultPaper}{HTML}{F0F8F4}
\hypersetup{colorlinks,linkcolor=BanglaGreen,urlcolor=BanglaGreen,
  citecolor=BanglaGreen,pdftitle={BOOKbangla 1.0.1: User guide},pdfauthor={Adib Hasan}}
\setlength{\parindent}{0pt}
\setlength{\parskip}{4pt plus 1pt}
\setlength{\emergencystretch}{2em}
\setsecheadstyle{\Large\sffamily\bfseries\color{BanglaGreen}\raggedright}
\setsubsecheadstyle{\normalsize\sffamily\bfseries\color{BanglaInk}\raggedright}
\setbeforesecskip{-9pt}
\setaftersecskip{6pt}
\setbeforesubsecskip{-8pt}
\setaftersubsecskip{3pt}
\setsecnumdepth{section}
\setsecnumformat{\csname the#1\endcsname\hspace{.6em}}
\AddToHook{cmd/section/before}{\Needspace{5\baselineskip}}
\AddToHook{cmd/subsection/before}{\Needspace{6\baselineskip}}
% English navigation, with ordinary Bangla numbering inside the examples.
\renewcommand*\thesection{\number\value{section}}
\renewcommand*\thepage{\number\value{page}}
\makepagestyle{manual}
\makeoddhead{manual}{\small\sffamily\color{BanglaGreen}BOOKbangla}{}%
  {\small\sffamily\color{BanglaInk}Version 1.0.1\quad User guide}
\makeoddfoot{manual}{}{\small\thepage}{}
\makeheadrule{manual}{\textwidth}{.3pt}
\pagestyle{manual}

% Execute the displayed source to produce each adjacent result panel.
\tcbset{manualbox/.style={enhanced,boxrule=.4pt,arc=2mm,
  colframe=BanglaGreen!35!white,colback=CodePaper,
  colbacktitle=BanglaGreen!9!white,coltitle=BanglaGreen,
  fonttitle=\small\sffamily\bfseries,left=3mm,right=3mm,top=2mm,bottom=2mm,
  before skip=5pt,after skip=5pt}}
\setminted{style=default,breaklines=true,breakanywhere=true}
\newminted[code]{latex}{fontsize=\small,samepage}
\newcommand*\manualcodetitle{Code}
\newcommand*\manualresulttitle{Result}
\newcounter{manualexample}
\newcommand*\manualexamplefile{\jobname-example-\number\value{manualexample}.tex}
\newenvironment{examplepair}
  {\VerbatimEnvironment\stepcounter{manualexample}%
   \begin{VerbatimOut}{\manualexamplefile}}
  {\end{VerbatimOut}%
   \begin{tcbraster}[raster columns=2,raster column skip=3mm,
     raster equal height=rows,raster valign=top]
   \begin{tcolorbox}[manualbox,title=\manualcodetitle]
     \inputminted[fontsize=\footnotesize]{latex}{\manualexamplefile}
   \end{tcolorbox}
   \begin{tcolorbox}[manualbox,title=\manualresulttitle,colback=ResultPaper]
     \input{\manualexamplefile}
   \end{tcolorbox}
   \end{tcbraster}}

% Chapter previews use the actual memoir heading at book width, reduced to
% fit the result panel. Only the surrounding page whitespace is omitted.
\makeatletter
\newcommand*\manualchapterpreview[1]{%
  \resizebox{\linewidth}{!}{\begin{minipage}{5in}
    \setlength{\beforechapskip}{0pt}\setlength{\afterchapskip}{0pt}%
    \m@m@Andtrue\@makechapterhead{#1}%
  \end{minipage}}}
\makeatother
\begin{document}
\selectlanguage{bengali}
\setcounter{chapter}{1}
\hypertarget{guide-english}{}
\phantomsection\label{manual:english}
\pdfbookmark[0]{English user guide}{manual-english}
\thispagestyle{empty}
\begin{center}
\includegraphics[width=.66\linewidth]{logo.png}\\[2pt]
{\Large\sffamily\color{BanglaGreen}Writing Bangla \& Assamese books with memoir}\\[5pt]
{Adib Hasan}\\[2pt]
{\small\href{https://groups.google.com/g/latexbangla}{\texttt{latexbangla@googlegroups.com}}}\\[3pt]
{24 October 2026\quad\textcolor{BanglaRed}{\textbullet}\quad Version 1.0.1}
\end{center}
Follow this guide in \hyperref[manual:bangla]{বাংলা} or
\hyperref[manual:assamese]{অসমীয়া}.

BOOKbangla helps you turn a Bangla or Assamese manuscript into a book without having to
design every page from scratch. Familiar \LaTeX{}
commands let you write chapters, add worked examples and problems, and refer
to them with consistent styling and chapter-based numbering.

The design comes from \emph{কম্বিনেটরিকসে হাতেখড়ি}\footnote{Adib Hasan,
Joydeep Saha and Ahmed Zawad Chowdhury, \emph{Introduction to Combinatorics},
Vol.~1, Dimik Prokashoni, 2019.} and can be adapted using
memoir's layout commands. \LaTeX{}bangla handles the selected language's text and numbering,
while you choose the text and mathematics fonts.

BOOKbangla is free to use under LPPL 1.3c or later. If you use it to publish
a book, I would appreciate an acknowledgement of BOOKbangla and its author,
Adib Hasan.

\section{Quickstart}
Use \hologo{XeLaTeX} or \hologo{LuaLaTeX} with an updated \TeX{} Live 2024 or later
and \LaTeX{}bangla 2.1 or newer. Install the
contents of the package directory in your local TeX tree, keeping its
\texttt{code/} directory, or copy those files beside your main document.
The following is a complete starting point:
\begin{tcolorbox}[manualbox,title=\manualcodetitle]
\begin{code}
\documentclass[12pt,twoside]{bookbangla}
\setbanglamainfont{Kalpurush}
\title{আমার বই}
\author{লেখকের নাম}
\date{}
\begin{document}
\frontmatter
\maketitle
\tableofcontents*
\mainmatter
\chapter{শুরু}
এখান থেকে বইয়ের লেখা শুরু।
\backmatter
\end{document}
\end{code}
\end{tcolorbox}
Choose an installed Bangla font or load a font file using \LaTeX{}bangla's setters.
English fonts remain independent. For example, this optional line selects
an English font:
\begin{tcolorbox}[manualbox,title=\manualcodetitle]
\begin{code}
\setmainfont{TeX Gyre Pagella}
\end{code}
\end{tcolorbox}
Use \verb|\setbanglasansfont| and \verb|\setbanglamonofont| for separate
sans and mono families in the selected language. No quotation font is required.

For an Assamese book, add \texttt{language=assamese} to the \verb|\documentclass| options.
Bangla is the default. The font commands stay the same; choose a font that
includes ৰ and ৱ. Headings, numbering and reference names follow your choice.

The example in \texttt{examples/bangla/main.tex} includes a title page, contents,
a complete combinatorics chapter, exercises and bibliography.
Compile with LaTeX, BibTeX, then LaTeX twice more.

\section{Title pages and chapter design}
The title page uses memoir's ordinary \verb|\maketitle| and the
\verb|\title|, \verb|\author| and \verb|\date| metadata. Put
\verb|\maketitle| inside a \texttt{titlingpage} environment for a separate
page. Nothing is inserted automatically. Write your own \texttt{titlingpage}
and omit \verb|\maketitle| for a completely custom design; copyright text
also belongs to the document.

Choose a chapter style with memoir's \verb|\chapterstyle| command.
\texttt{classic} is the default. All four are suitable for black-and-white printing;
other memoir chapter styles remain available too. Use
\verb|\chapterstyle{default}| for memoir's normal chapter heading. This
changes chapter headings, not the page layout or section styling. See
\texttt{examples/chapter-styles.pdf} for full-page samples. The previews below
show the same chapter title in each design, reduced to fit the panels.
\begingroup
\let\chapter\manualchapterpreview
\begin{examplepair}
\chapterstyle{classic}
\chapter{গণনার শুরু}
\end{examplepair}
\begin{examplepair}
\chapterstyle{modern}
\chapter{গণনার শুরু}
\end{examplepair}
\begin{examplepair}
\chapterstyle{minimal}
\chapter{গণনার শুরু}
\end{examplepair}
\begin{examplepair}
\chapterstyle{ruled}
\chapter{গণনার শুরু}
\end{examplepair}
\endgroup

\subsection{Make the design your own}

Use \verb|\bookbanglasetup{...}| to set the following options.

\begin{center}
\begin{tabular}{ll}
\toprule
Setting & Meaning\\
\midrule
\texttt{accent-color} & Chapter accents (default dark gray)\\
\texttt{box-frame-color} & Box borders and titles (default black)\\
\texttt{box-background-color} & Box backgrounds (default white)\\
\texttt{quote-font} & Quotation font commands (default Chorus/italic)\\
\bottomrule
\end{tabular}
\end{center}
These are BOOKbangla's own settings. For other adjustments, use memoir after
loading the class. Later explicit changes take precedence. For example:
\begin{tcolorbox}[manualbox,title=\manualcodetitle]
\begin{code}
\chapterstyle{modern} % Replace BOOKbangla's chapter style.
\pagestyle{plain}
\setlength{\parindent}{1em}
\setlength{\parskip}{0pt}
\end{code}
\end{tcolorbox}
Return to the default chapter design with \verb|\chapterstyle{classic}|. The section
heading's \S{} marker is decoration, so references contain only the section
number. To remove it, use memoir's formatter:
\begin{tcolorbox}[manualbox,title=\manualcodetitle]
\begin{code}
\setsecnumformat{\csname the#1\endcsname\quad}
\end{code}
\end{tcolorbox}
The default trimmed page is 9.5 by 6.8 inches on A4 stock. Change trim size,
type block and margins using memoir, then call \verb|\checkandfixthelayout|.
Smaller paper options, such as \verb|a5paper|, use memoir's page dimensions
and margins instead.
For an A5 book without trimming:
\begin{tcolorbox}[manualbox,title=\manualcodetitle]
\begin{code}
\setstocksize{210mm}{148mm}
\settrimmedsize{210mm}{148mm}{*}
\settrims{0pt}{0pt}
\setlrmarginsandblock{20mm}{18mm}{*}
\setulmarginsandblock{22mm}{25mm}{*}
\checkandfixthelayout
\end{code}
\end{tcolorbox}
Running headers use the book title on even pages and the section mark on odd
pages. Customize them with memoir's page-style commands. The
\href{https://ctan.org/pkg/memoir}{memoir manual} describes these interfaces.

\section{Examples, statements and exercises}
An \texttt{example} is a numbered, breakable box. Its numbering restarts in
each chapter, and both ordinary labels and the \texttt{label} option
work:
\begin{examplepair}
\begin{example}[label=ex:weekdays]
আটজন বন্ধুর মধ্যে অন্তত দুজনের জন্ম একই বারে কেন?
\end{example}
\begin{solution}
সপ্তাহে সাতটি বার আছে।
\end{solution}
দেখো \autoref{ex:weekdays}।
\end{examplepair}
\texttt{solution} uses proof-style presentation and an end mark; give it an
optional heading with \verb|\begin{solution}[অন্য সমাধান]|. The ordinary
\texttt{proof} environment remains available, including \verb|\qedhere|.
Use \texttt{remark} for an unnumbered মন্তব্য.

\texttt{reminder} supplies the মনে রাখবে box with an optional replacement
title, and \texttt{diybox} supplies the নিজে করো box. For a formal statement,
\texttt{theorem} uses amsthm's ordinary plain style: a bold heading and italic
body, with its own chapter-based numbering. An optional argument adds the
theorem's name; ordinary labels and references work as usual:
\begin{examplepair}
\begin{theorem}[পায়রার খোপ নীতি]
\label{thm:pigeonhole}
তিনটি বস্তু দুটি খোপে রাখলে
কোনো খোপে অন্তত দুটি বস্তু থাকে।
\end{theorem}
\begin{proof}
প্রতি খোপে সর্বোচ্চ একটি বস্তু
রাখলে দুটি বস্তু রাখা যায়।
\end{proof}
দেখো \autoref{thm:pigeonhole}।
\end{examplepair}
\begin{examplepair}
\begin{reminder}
বস্তু ও খোপ আগে চিহ্নিত করো।
\end{reminder}
\begin{diybox}
চারটি বস্তু নিয়ে চেষ্টা করো।
\end{diybox}
\end{examplepair}
Define additional statement environments with \verb|\newtheorem| when needed,
choosing whether they share the theorem counter or have their own.

Each \texttt{exercise} is one problem, labelled সমস্যা. Use an ordinary
section to group problems:
\begin{tcolorbox}[manualbox,title=\manualcodetitle]
\begin{code}
\section{অনুশীলনী}
\end{code}
\end{tcolorbox}
\begin{examplepair}
\begin{exercise}[
  hint={মাসগুলোকে খোপ ধরো;
        $13>12$।}]
\label{ex:months}
তেরোজন মানুষের মধ্যে অন্তত দুজনের জন্ম একই মাসে, দেখাও।
\end{exercise}
\begin{exercise}
চারটি পূর্ণসংখ্যার মধ্যে দুটির
$3$ দিয়ে ভাগশেষ সমান, দেখাও।
\end{exercise}
দেখো \autoref{ex:months}।
\end{examplepair}
A supplied hint appears after the problem under ইঙ্গিত:. Without a hint,
only the problem is printed. Exercise numbering is separate from examples,
continues across sections, and restarts at each chapter. Both support
\verb|\label|, \verb|\ref| and \verb|\autoref|. The \texttt{cleveref} package also
uses translated names for these environments. Appendices follow memoir's
chapter numbering, localized by \LaTeX{}bangla. To number problems within
sections instead, use memoir's counter command in the preamble:
\begin{tcolorbox}[manualbox,title=\manualcodetitle]
\begin{code}
\counterwithin{exercise}{section}
\end{code}
\end{tcolorbox}
Use \verb|\counterwithout{exercise}{chapter}| for a continuous sequence across
the book. The same commands apply to the \texttt{example} counter.

\section{Quotations, lists and mathematics}
\begin{examplepair}
\begin{chapquote}[1em]{লেখক}
ছোট প্রশ্ন থেকেই বড় ধারণার জন্ম হয়।
\end{chapquote}
\end{examplepair}
The optional argument sets both side indents; the author’s name follows
the quotation. English uses TeX Gyre Chorus at 105\% scale by default; Bangla and Assamese
use the italic form of your chosen font. To choose another English quotation font:
\begin{tcolorbox}[manualbox,title=\manualcodetitle]
\begin{code}
\newfontfamily\quoteface{TeX Gyre Cursor}
\bookbanglasetup{quote-font={\quoteface\itshape}}
\end{code}
\end{tcolorbox}
\LaTeX{}bangla handles Bangla or Assamese within the quotation. A decorative
opening can use a drop cap:
\begin{examplepair}
\firstword{আ}{মরা} পরিচিত একটি
ধারণা দিয়ে শুরু করি। ছোট প্রশ্ন
থেকে অনেক বড় ধারণা জন্ম নেয়।
\end{examplepair}

\texttt{smartlist} uses a diamond marker by default. Its options are
\texttt{enumitem} options, so labels, margins and spacing can be changed locally:
\begin{examplepair}
\begin{smartlist}
\item বস্তুগুলো চিহ্নিত করো।
\item খোপগুলো গুনে দেখো।
\end{smartlist}
\begin{smartlist}[
  label=$\Rightarrow$,labelsep=.5em,
  leftmargin=2em,itemsep=0pt]
\item নিজের যুক্তি লিখে ফেলো।
\end{smartlist}
\end{examplepair}
Ordinary \texttt{itemize} and \texttt{enumerate} are unchanged.

AMS mathematics and proof support are loaded. Use \verb|\eqref| for equation
references with parentheses. Multiline displays may break across pages; use
\verb|\allowdisplaybreaks[0]| to prevent that. Graphics and table footnotes
are supported. For other features, use memoir or packages such as
\texttt{longtable}, \texttt{multicol} and \texttt{placeins}.

\section{Bibliography and package options}
Choose your bibliography package, style and resources in the document. The
class does not load \texttt{natbib} or \texttt{biblatex}; the example book uses
\texttt{biblatex} with Biber. It loads
\texttt{hyperref}; change link appearance with \verb|\hypersetup|. To change
\LaTeX{}bangla's numbering option, pass it before loading the class:
\begin{tcolorbox}[manualbox,title=\manualcodetitle]
\begin{code}
\PassOptionsToPackage{changecounternumbering=false}{latexbangla}
\documentclass{bookbangla}
\end{code}
\end{tcolorbox}

\smallskip
Copyright 2019, 2026 Adib Hasan. Maintained by Adib Hasan.
See \texttt{README} for the list of files.

\clearpage
% Keep page numbers continuous, but restart the translated guide's sections
% and demonstrations. Distinct labels and PDF anchors keep references local.
\renewcommand*\thesection{\arabic{section}}
\renewcommand*\thepage{\arabic{page}}
\renewcommand*\theHsection{bangla.\number\value{section}}
\renewcommand*\theHexample{bangla.\number\value{example}}
\renewcommand*\theHexercise{bangla.\number\value{exercise}}
\renewcommand*\theHtheorem{bangla.\number\value{theorem}}
\setcounter{section}{0}
\setcounter{example}{0}
\setcounter{exercise}{0}
\setcounter{theorem}{0}
\renewcommand*\manualcodetitle{কোড}
\renewcommand*\manualresulttitle{ফলাফল}
\makeoddhead{manual}{\small\sffamily\color{BanglaGreen}BOOKbangla}{}%
  {\small\sffamily\color{BanglaInk}সংস্করণ ১.০.১\quad ব্যবহার নির্দেশিকা}
\hypertarget{guide-bangla}{}
\phantomsection\label{manual:bangla}
\pdfbookmark[0]{বাংলা ব্যবহার নির্দেশিকা}{manual-bangla}
\thispagestyle{empty}
\begin{center}
\includegraphics[width=.66\linewidth]{logo.png}\\[2pt]
{\Large\sffamily\color{BanglaGreen}memoir দিয়ে বাংলা ও অসমীয়া বই লেখা}\\[5pt]
{আদীব হাসান}\\[2pt]
{\small\href{https://groups.google.com/g/latexbangla}{\texttt{latexbangla@googlegroups.com}}}\\[3pt]
{২৪ অক্টোবর ২০২৬\quad\textcolor{BanglaRed}{\textbullet}\quad সংস্করণ ১.০.১}
\end{center}
এই নির্দেশিকা \hyperref[manual:english]{ইংরেজি} ও
\hyperref[manual:assamese]{অসমীয়াতেও} পড়তে পারেন।

BOOKbangla দিয়ে বাংলা বা অসমীয়া পাণ্ডুলিপি থেকে বই তৈরি করতে প্রতিটি পৃষ্ঠার নকশা
গোড়া থেকে করতে হয় না। পরিচিত \LaTeX-এর কমান্ড দিয়েই অধ্যায়
লিখতে, সমাধানসহ উদাহরণ ও সমস্যা যোগ করতে এবং সেগুলোর রেফারেন্স দিতে পারেন।
এগুলোর সাজানোর রীতি ও অধ্যায়ভিত্তিক নম্বর ক্লাস নিজেই সামলায়।

নকশাটি এসেছে \emph{কম্বিনেটরিকসে হাতেখড়ি}\footnote{আদীব হাসান,
জয়দীপ সাহা ও আহমেদ জাওয়াদ চৌধুরী, \emph{কম্বিনেটরিকসে হাতেখড়ি — ১ম খণ্ড},
দ্বিমিক প্রকাশনী, ২০১৯।} থেকে; memoir-এর কমান্ড দিয়ে
পৃষ্ঠার বিন্যাস নিজের মতো বদলে নিতে পারেন। নির্বাচিত ভাষার লেখা ও নম্বর সামলায়
\LaTeX{}bangla, আর লেখার ও গণিতের ফন্ট বেছে নেন আপনি।

BOOKbangla LPPL 1.3c বা পরবর্তী সংস্করণের অধীনে বিনামূল্যে ব্যবহার করতে পারেন।
এটি দিয়ে বই প্রকাশ করলে বইয়ে BOOKbangla ও এর লেখক আদীব হাসানের নাম
উল্লেখ করলে খুশি হব।

\section{দ্রুত আরম্ভ}
হালনাগাদ \TeX{} Live 2024 বা পরবর্তী সংস্করণে \hologo{XeLaTeX} বা \hologo{LuaLaTeX}
ব্যবহার করুন। \LaTeX{}bangla ২.১ বা পরবর্তী
সংস্করণ ইনস্টল করা থাকতে হবে। প্যাকেজের ফাইলগুলো আপনার কম্পিউটারের স্থানীয়
TeX ট্রিতে ইনস্টল করুন, অথবা মূল ডকুমেন্টের পাশে রাখুন। দুই ক্ষেত্রেই
\texttt{code/} ফোল্ডারের বিন্যাস অক্ষত রাখুন।
শুরু করার জন্য নিচের সম্পূর্ণ ডকুমেন্টটি ব্যবহার করতে পারেন:
\begin{tcolorbox}[manualbox,title=\manualcodetitle]
\begin{code}
\documentclass[12pt,twoside]{bookbangla}
\setbanglamainfont{Kalpurush}
\title{আমার বই}
\author{লেখকের নাম}
\date{}
\begin{document}
\frontmatter
\maketitle
\tableofcontents*
\mainmatter
\chapter{শুরু}
এখান থেকে বইয়ের লেখা শুরু।
\backmatter
\end{document}
\end{code}
\end{tcolorbox}
\LaTeX{}bangla-এর ফন্ট নির্ধারণের কমান্ড দিয়ে ইনস্টল করা বাংলা ফন্ট বেছে নিতে
বা সরাসরি ফন্ট ফাইল লোড করতে পারেন। ইংরেজি ফন্ট আলাদাভাবে বেছে নেওয়া যায়।
যেমন, ইংরেজি ফন্ট বদলাতে চাইলে এই লাইনটি যোগ করতে পারেন:
\begin{tcolorbox}[manualbox,title=\manualcodetitle]
\begin{code}
\setmainfont{TeX Gyre Pagella}
\end{code}
\end{tcolorbox}
\verb|\setbanglasansfont| ও \verb|\setbanglamonofont| দিয়ে নির্বাচিত ভাষার
sans ও mono ফন্ট আলাদা করে বেছে নিতে পারেন। উদ্ধৃতির জন্য আলাদা ফন্ট দেওয়া জরুরি নয়।

অসমীয়া বইয়ের জন্য ক্লাসের অপশন দিন \texttt{language=assamese}।
ডিফল্ট ভাষা বাংলা। ফন্টের কমান্ডগুলো একই থাকে; ৰ ও ৱ আছে এমন ফন্ট বেছে নিন।
শিরোনাম, নম্বর ও রেফারেন্সের নাম নির্বাচিত ভাষা অনুযায়ী আসে।

সঙ্গে দেওয়া \texttt{examples/bangla/main.tex}-এ শিরোনাম পৃষ্ঠা, সূচিপত্র,
কম্বিনেটরিকসের একটি সম্পূর্ণ অধ্যায়, অনুশীলনী ও গ্রন্থসূত্র আছে। প্রথমে LaTeX,
তারপর BibTeX, এরপর আরও দুবার LaTeX চালিয়ে কম্পাইল করুন।

\section{শিরোনাম পৃষ্ঠা ও অধ্যায়ের নকশা}
শিরোনাম পৃষ্ঠাটি memoir-এর পরিচিত \verb|\maketitle| কমান্ড দিয়ে তৈরি হয়;
বইয়ের নাম, লেখক ও তারিখ আসে \verb|\title|, \verb|\author| ও \verb|\date|
থেকে। শিরোনামের জন্য আলাদা পৃষ্ঠা রাখতে \texttt{titlingpage} এনভায়রনমেন্টের ভেতরে
\verb|\maketitle| লিখুন। ক্লাস নিজে থেকে শিরোনাম পৃষ্ঠা যোগ করে না। সম্পূর্ণ নিজের
নকশা চাইলে \verb|\maketitle| বাদ দিয়ে \texttt{titlingpage}-এর ভেতরে সেটি
লিখুন। কপিরাইটের লেখাও আপনার ডকুমেন্টেই রাখবেন।

memoir-এর \verb|\chapterstyle| কমান্ড দিয়ে অধ্যায়ের নকশা বেছে নিন।
ডিফল্ট নকশা হলো \texttt{classic}। চারটি নকশাই সাদাকালো ছাপার উপযোগী;
memoir-এর অন্য নকশাগুলোও ব্যবহার করতে পারেন। memoir-এর সাধারণ অধ্যায়ের
শিরোনাম পেতে \verb|\chapterstyle{default}| লিখুন। এতে শুধু অধ্যায়ের শিরোনাম
বদলায়; পৃষ্ঠার বিন্যাস বা পরিচ্ছেদের নকশা বদলায় না। পুরো পৃষ্ঠায় নকশাগুলো
দেখতে \texttt{examples/chapter-styles.pdf} খুলুন। নিচে একই অধ্যায়ের শিরোনাম
চারটি নকশায় দেখানো হয়েছে; বক্সে ধরানোর জন্য আকার ছোট করা হয়েছে।
\begingroup
\let\chapter\manualchapterpreview
\begin{examplepair}
\chapterstyle{classic}
\chapter{গণনার শুরু}
\end{examplepair}
\begin{examplepair}
\chapterstyle{modern}
\chapter{গণনার শুরু}
\end{examplepair}
\begin{examplepair}
\chapterstyle{minimal}
\chapter{গণনার শুরু}
\end{examplepair}
\begin{examplepair}
\chapterstyle{ruled}
\chapter{গণনার শুরু}
\end{examplepair}
\endgroup

\subsection{নিজের মতো করে সাজানো}

নিচের সেটিংগুলো বদলাতে \verb|\bookbanglasetup{...}| ব্যবহার করুন।

\begin{center}
\begin{tabular}{ll}
\toprule
সেটিং & কাজ\\
\midrule
\texttt{accent-color} & অধ্যায়ের নকশার রং (ডিফল্ট গাঢ় ধূসর)\\
\texttt{box-frame-color} & বক্সের সীমানা ও শিরোনাম (ডিফল্ট কালো)\\
\texttt{box-background-color} & বক্সের পটভূমি (ডিফল্ট সাদা)\\
\texttt{quote-font} & উদ্ধৃতির ফন্টের কমান্ড (ডিফল্ট Chorus/তির্যক)\\
\bottomrule
\end{tabular}
\end{center}
এগুলো BOOKbangla-এর নিজস্ব সেটিং। অন্য কিছু বদলাতে ক্লাস লোড করার পরে
memoir-এর কমান্ড ব্যবহার করুন। পরে দেওয়া নির্দেশই কার্যকর হবে। যেমন:
\begin{tcolorbox}[manualbox,title=\manualcodetitle]
\begin{code}
\chapterstyle{modern} % অধ্যায়ের নকশা বদলান।
\pagestyle{plain}
\setlength{\parindent}{1em}
\setlength{\parskip}{0pt}
\end{code}
\end{tcolorbox}
ডিফল্ট নকশায় ফিরতে \verb|\chapterstyle{classic}| লিখুন। পরিচ্ছেদের শিরোনামের
\S{} চিহ্নটি শুধু সাজানোর জন্য; রেফারেন্সে কেবল পরিচ্ছেদের নম্বর থাকে।
চিহ্নটি বাদ দিতে memoir-এর নম্বর সাজানোর কমান্ড ব্যবহার করুন:
\begin{tcolorbox}[manualbox,title=\manualcodetitle]
\begin{code}
\setsecnumformat{\csname the#1\endcsname\quad}
\end{code}
\end{tcolorbox}
ডিফল্ট বিন্যাসে A4 কাগজ কেটে তৈরি পৃষ্ঠার মাপ ৯.৫ বাই ৬.৮ ইঞ্চি। memoir দিয়ে
কাটার পরের মাপ, লেখার এলাকার মাপ ও মার্জিন বদলে \verb|\checkandfixthelayout|
চালান। \verb|a5paper|-এর মতো ছোট কাগজের অপশন দিলে memoir-এর পৃষ্ঠার মাপ
ও মার্জিন ব্যবহৃত হয়। কাগজ না কেটে সরাসরি A5 মাপের বইয়ের জন্য:
\begin{tcolorbox}[manualbox,title=\manualcodetitle]
\begin{code}
\setstocksize{210mm}{148mm}
\settrimmedsize{210mm}{148mm}{*}
\settrims{0pt}{0pt}
\setlrmarginsandblock{20mm}{18mm}{*}
\setulmarginsandblock{22mm}{25mm}{*}
\checkandfixthelayout
\end{code}
\end{tcolorbox}
জোড় পৃষ্ঠার ওপরে বইয়ের নাম এবং বিজোড় পৃষ্ঠার ওপরে পরিচ্ছেদের নম্বর ও
শিরোনাম থাকে। memoir-এর পৃষ্ঠাশৈলীর কমান্ড দিয়ে এগুলো বদলাতে পারেন।
বিস্তারিত পাবেন \href{https://ctan.org/pkg/memoir}{memoir-এর নির্দেশিকায়}।

\section{উদাহরণ, উপপাদ্য ও অনুশীলনী}
\texttt{example} হলো নম্বরযুক্ত উদাহরণের বক্স; প্রয়োজনে এটি একাধিক পৃষ্ঠায়
ছড়িয়ে যেতে পারে। প্রতি অধ্যায়ে এর নম্বর নতুন করে শুরু হয়। সাধারণ
\verb|\label| কমান্ড ও \texttt{label} অপশন দুটোই ব্যবহার করা যায়:
\begin{examplepair}
\begin{example}[label=bn:weekdays]
আটজন বন্ধুর মধ্যে অন্তত দুজনের জন্ম একই বারে কেন?
\end{example}
\begin{solution}
সপ্তাহে সাতটি বার আছে।
\end{solution}
দেখো \autoref{bn:weekdays}।
\end{examplepair}
\texttt{solution} প্রমাণের মতো করে সমাধান সাজায় এবং শেষে সমাপ্তিচিহ্ন দেয়।

\verb|\begin{solution}[অন্য সমাধান]| দিয়ে শিরোনাম বদলাতে পারেন। পরিচিত
\texttt{proof} এনভায়রনমেন্ট ও \verb|\qedhere| কমান্ডও ব্যবহার করা যায়।
নম্বর ছাড়া মন্তব্য লিখতে \texttt{remark} ব্যবহার করুন।

\texttt{reminder} দিয়ে মনে রাখবে বক্স তৈরি হয়; ঐচ্ছিক আর্গুমেন্ট দিয়ে এর
শিরোনাম বদলাতে পারেন। \texttt{diybox} দেয় নিজে করো বক্স। উপপাদ্যের জন্য
\texttt{theorem} ব্যবহার করুন। এতে amsthm-এর সাধারণ \texttt{plain} শৈলী
অনুযায়ী গাঢ় শিরোনাম ও তির্যক লেখা থাকে, আর নম্বর চলে অধ্যায়ভিত্তিক আলাদা
ধারায়। ঐচ্ছিক আর্গুমেন্টে উপপাদ্যের নাম যোগ করা যায়; পরিচিত লেবেল ও
রেফারেন্সের কমান্ডও কাজ করে:
\begin{examplepair}
\begin{theorem}[পায়রার খোপ নীতি]
\label{bn:theorem}
তিনটি বস্তু দুটি খোপে রাখলে
কোনো খোপে অন্তত দুটি বস্তু থাকে।
\end{theorem}
\begin{proof}
প্রতি খোপে সর্বোচ্চ একটি বস্তু
রাখলে দুটি বস্তু রাখা যায়।
\end{proof}
দেখো \autoref{bn:theorem}।
\end{examplepair}
\begin{examplepair}
\begin{reminder}
বস্তু ও খোপ আগে চিহ্নিত করো।
\end{reminder}
\begin{diybox}
চারটি বস্তু নিয়ে চেষ্টা করো।
\end{diybox}
\end{examplepair}
আরও উপপাদ্য বা অনুরূপ এনভায়রনমেন্ট দরকার হলে \verb|\newtheorem| দিয়ে তৈরি
করুন। সেগুলোর নম্বর উপপাদ্যের সঙ্গে একই ধারায় চলবে, নাকি আলাদা হবে, তা
আপনিই বেছে নিতে পারেন।

প্রতিটি \texttt{exercise}-এ একটি সমস্যা লিখবেন; এর শিরোনাম হয় সমস্যা।
সমস্যাগুলো একসঙ্গে রাখতে সাধারণ পরিচ্ছেদ ব্যবহার করুন:
\begin{tcolorbox}[manualbox,title=\manualcodetitle]
\begin{code}
\section{অনুশীলনী}
\end{code}
\end{tcolorbox}
\begin{examplepair}
\begin{exercise}[
  hint={মাসগুলোকে খোপ ধরো;
        $13>12$।}]
\label{bn:months}
তেরোজন মানুষের মধ্যে অন্তত দুজনের জন্ম একই মাসে, দেখাও।
\end{exercise}
\begin{exercise}
চারটি পূর্ণসংখ্যার মধ্যে দুটির
$3$ দিয়ে ভাগশেষ সমান, দেখাও।
\end{exercise}
দেখো \autoref{bn:months}।
\end{examplepair}
ইঙ্গিত দিলে সেটি সমস্যার পরে ইঙ্গিত: শিরোনামে দেখা যায়। না দিলে শুধু
সমস্যাটিই ছাপা হয়। সমস্যার নম্বর উদাহরণের নম্বর থেকে আলাদা; পরিচ্ছেদ বদলালে
তা চলতে থাকে, নতুন অধ্যায়ে আবার শুরু হয়। দুটোতেই \verb|\label|,
\verb|\ref| ও \verb|\autoref| ব্যবহার করা যায়। \texttt{cleveref}-ও এই এনভায়রনমেন্টগুলোর
অনূদিত নাম ব্যবহার করে। পরিশিষ্টে memoir-এর অধ্যায়ের নম্বর ব্যবহৃত হয়, যা \LaTeX{}bangla বাংলায় দেখায়। প্রতি পরিচ্ছেদে সমস্যার
নম্বর নতুন করে শুরু করতে চাইলে প্রিঅ্যাম্বলে লিখুন:
\begin{tcolorbox}[manualbox,title=\manualcodetitle]
\begin{code}
\counterwithin{exercise}{section}
\end{code}
\end{tcolorbox}
পুরো বইয়ে সমস্যার নম্বর একটানা চালাতে
\verb|\counterwithout{exercise}{chapter}| ব্যবহার করুন। \texttt{example}
কাউন্টারের ক্ষেত্রেও একই কমান্ডগুলো প্রযোজ্য।

\section{উদ্ধৃতি, তালিকা ও গণিত}
\begin{examplepair}
\begin{chapquote}[1em]{লেখক}
ছোট প্রশ্ন থেকেই বড় ধারণার জন্ম হয়।
\end{chapquote}
\end{examplepair}
ঐচ্ছিক আর্গুমেন্ট দিয়ে উদ্ধৃতির দুই পাশের ফাঁকা জায়গা নির্ধারণ করা যায়;
শেষে লেখকের নাম থাকে। ডিফল্টভাবে ইংরেজিতে ১০৫\% স্কেলে TeX Gyre Chorus এবং বাংলা ও
অসমীয়ায় আপনার বেছে নেওয়া ফন্টের তির্যক রূপ ব্যবহৃত হয়। ইংরেজির ফন্ট বদলাতে:
\begin{tcolorbox}[manualbox,title=\manualcodetitle]
\begin{code}
\newfontfamily\quoteface{TeX Gyre Cursor}
\bookbanglasetup{quote-font={\quoteface\itshape}}
\end{code}
\end{tcolorbox}
উদ্ধৃতির ভেতরের বাংলা বা অসমীয়া লেখা \LaTeX{}bangla সামলায়। অনুচ্ছেদের
শুরুতে প্রথম অক্ষরটি বড় করে সাজাতে পারেন:
\begin{examplepair}
\firstword{আ}{মরা} পরিচিত একটি
ধারণা দিয়ে শুরু করি। ছোট প্রশ্ন
থেকে অনেক বড় ধারণা জন্ম নেয়।
\end{examplepair}

\texttt{smartlist}-এ ডিফল্টভাবে হীরকাকৃতির চিহ্ন থাকে। এতে \texttt{enumitem}-এর
অপশনগুলো ব্যবহার করে নির্দিষ্ট তালিকার চিহ্ন, মার্জিন ও ব্যবধান বদলাতে পারেন:
\begin{examplepair}
\begin{smartlist}
\item বস্তুগুলো চিহ্নিত করো।
\item খোপগুলো গুনে দেখো।
\end{smartlist}
\begin{smartlist}[
  label=$\Rightarrow$,labelsep=.5em,
  leftmargin=2em,itemsep=0pt]
\item নিজের যুক্তি লিখে ফেলো।
\end{smartlist}
\end{examplepair}
পরিচিত \texttt{itemize} ও \texttt{enumerate} আগের মতোই কাজ করে।

AMS-এর গণিত ও প্রমাণের সুবিধাগুলো লোড করা আছে। বন্ধনীসহ সমীকরণের রেফারেন্স
দিতে \verb|\eqref| ব্যবহার করুন। একাধিক লাইনের সমীকরণ পৃষ্ঠা পেরিয়ে চলতে
পারে; তা বন্ধ করতে \verb|\allowdisplaybreaks[0]| লিখুন। ছবি যোগ করা ও
সারণিতে পাদটীকা দেওয়াও যায়। অন্য সুবিধার জন্য memoir অথবা
\texttt{longtable}, \texttt{multicol}, \texttt{placeins}-এর মতো প্যাকেজ ব্যবহার করুন।

\section{গ্রন্থসূত্র ও প্যাকেজের অপশন}
গ্রন্থসূত্রের প্যাকেজ, বিন্যাস ও তথ্যের ফাইল আপনার ডকুমেন্টে বেছে নিন। ক্লাস
\texttt{natbib} বা \texttt{biblatex} লোড করে না; উদাহরণের বইয়ে
\texttt{biblatex} ও Biber ব্যবহার করা হয়েছে। ক্লাস \texttt{hyperref} লোড করে;
লিংকের চেহারা বদলাতে \verb|\hypersetup| ব্যবহার করুন। \LaTeX{}bangla-এর
নম্বরের অপশন বদলাতে ক্লাস লোড করার আগে সেটি দিন:
\begin{tcolorbox}[manualbox,title=\manualcodetitle]
\begin{code}
\PassOptionsToPackage{changecounternumbering=false}{latexbangla}
\documentclass{bookbangla}
\end{code}
\end{tcolorbox}

\smallskip
কপিরাইট ২০১৯, ২০২৬ আদীব হাসান। রক্ষণাবেক্ষণকারী আদীব হাসান।
ফাইলের তালিকা দেখুন \texttt{README}-তে।

\clearpage
\selectlanguage{assamese}
\renewcommand*\theHsection{assamese.\number\value{section}}
\renewcommand*\theHexample{assamese.\number\value{example}}
\renewcommand*\theHexercise{assamese.\number\value{exercise}}
\renewcommand*\theHtheorem{assamese.\number\value{theorem}}
\setcounter{section}{0}
\setcounter{example}{0}
\setcounter{exercise}{0}
\setcounter{theorem}{0}
\renewcommand*\manualcodetitle{কোড}
\renewcommand*\manualresulttitle{ফলাফল}
\makeoddhead{manual}{\small\sffamily\color{BanglaGreen}BOOKbangla}{}%
  {\small\sffamily\color{BanglaInk}সংস্কৰণ ১.০.১\quad ব্যৱহাৰ নিৰ্দেশিকা}
\hypertarget{guide-assamese}{}
\phantomsection\label{manual:assamese}
\pdfbookmark[0]{অসমীয়া ব্যৱহাৰ নিৰ্দেশিকা}{manual-assamese}
\thispagestyle{empty}
\begin{center}
\includegraphics[width=.66\linewidth]{logo.png}\\[2pt]
{\Large\sffamily\color{BanglaGreen}memoir-এৰে বাংলা আৰু অসমীয়া কিতাপ লিখক}\\[5pt]
{আদীব হাসান}\\[2pt]
{\small\href{https://groups.google.com/g/latexbangla}{\texttt{latexbangla@googlegroups.com}}}\\[3pt]
{২৪ অক্টোবৰ ২০২৬\quad\textcolor{BanglaRed}{\textbullet}\quad সংস্কৰণ ১.০.১}
\end{center}
এই নিৰ্দেশিকা \hyperref[manual:english]{ইংৰাজী} আৰু
\hyperref[manual:bangla]{বাংলাতো} পঢ়িব পাৰে।

BOOKbangla-ৰে বাংলা বা অসমীয়া পাণ্ডুলিপিৰ পৰা কিতাপ তৈয়াৰ কৰোঁতে প্ৰতিটো
পৃষ্ঠাৰ নক্সা আৰম্ভণিৰ পৰা কৰিব নালাগে। চিনাকি \LaTeX{} কমাণ্ডেৰেই অধ্যায়
লিখিব, সমাধানসহ উদাহৰণ আৰু সমস্যা যোগ কৰিব, আৰু সেইবোৰৰ উল্লেখ দিব পাৰে।
সেইবোৰৰ সজ্জা আৰু অধ্যায়ভিত্তিক নম্বৰ ক্লাছটোৱেই সামলায়।

নক্সাটো \emph{কম্বিনেটরিকসে হাতেখড়ি}\footnote{আদীব হাসান,
জয়দীপ সাহা আৰু আহমেদ জাওয়াদ চৌধুৰী, \emph{কম্বিনেটরিকসে হাতেখড়ি — ১ম খণ্ড},
দ্বিমিক প্ৰকাশনী, ২০১৯।} কিতাপৰ পৰা লোৱা হৈছে; memoir-ৰ কমাণ্ডেৰে
পৃষ্ঠাৰ বিন্যাস নিজৰ মতে সলনি কৰিব পাৰে। বাছনি কৰা ভাষাৰ লিখনি আৰু নম্বৰ
\LaTeX{}bangla-ই সামলায়, আৰু লিখনিৰ তথা গণিতৰ ফণ্ট আপুনি বাছি লয়।

BOOKbangla LPPL 1.3c বা পৰৱৰ্তী সংস্কৰণৰ অধীনত বিনামূলীয়াকৈ ব্যৱহাৰ কৰিব পাৰে।
ইয়াৰে কিতাপ প্ৰকাশ কৰিলে কিতাপখনত BOOKbangla আৰু ইয়াৰ লেখক আদীব হাসানৰ
নাম উল্লেখ কৰিলে আনন্দিত হ'ম।

\section{সোনকালে আৰম্ভ কৰক}
হালনাগাদ \TeX{} Live 2024 বা পৰৱৰ্তী সংস্কৰণত \hologo{XeLaTeX} অথবা \hologo{LuaLaTeX}
ব্যৱহাৰ কৰক। \LaTeX{}bangla ২.১ বা
পৰৱৰ্তী সংস্কৰণ ইনষ্টল কৰা থাকিব লাগিব। পেকেজৰ ফাইলবোৰ আপোনাৰ কম্পিউটাৰৰ
স্থানীয় TeX ট্ৰীত ইনষ্টল কৰক, অথবা মূল ডকুমেণ্টৰ কাষত ৰাখক। দুয়ো ক্ষেত্ৰতে
\texttt{code/} ফ'ল্ডাৰৰ গঠন অক্ষুণ্ণ ৰাখিব। আৰম্ভ কৰিবলৈ তলৰ সম্পূৰ্ণ
ডকুমেণ্টটো ব্যৱহাৰ কৰিব পাৰে:
\begin{tcolorbox}[manualbox,title=\manualcodetitle]
\begin{code}
\documentclass[language=assamese,12pt,twoside]{bookbangla}
\setbanglamainfont{Kalpurush}
\title{মোৰ কিতাপ}
\author{লেখকৰ নাম}
\date{}
\begin{document}
\frontmatter
\maketitle
\tableofcontents*
\mainmatter
\chapter{আৰম্ভণি}
ইয়াৰ পৰা কিতাপৰ লিখনি আৰম্ভ।
\backmatter
\end{document}
\end{code}
\end{tcolorbox}
\LaTeX{}bangla-ৰ ফণ্ট নিৰ্ধাৰণৰ কমাণ্ডেৰে ইনষ্টল কৰা ফণ্ট বাছি ল'ব পাৰে,
অথবা পোনপটীয়াকৈ ফণ্ট ফাইল লোড কৰিব পাৰে। ইংৰাজী ফণ্ট পৃথককৈ বাছিব পাৰে।
উদাহৰণস্বৰূপে, ইংৰাজী ফণ্ট সলনি কৰিবলৈ এই শাৰীটো যোগ কৰক:
\begin{tcolorbox}[manualbox,title=\manualcodetitle]
\begin{code}
\setmainfont{TeX Gyre Pagella}
\end{code}
\end{tcolorbox}
বাছনি কৰা ভাষাৰ sans আৰু mono পৰিয়ালৰ বাবে পৃথক ফণ্ট বিচাৰিলে
\verb|\setbanglasansfont| আৰু \verb|\setbanglamonofont| ব্যৱহাৰ কৰক।
উদ্ধৃতিৰ বাবে পৃথক ফণ্ট দিয়াৰ প্ৰয়োজন নাই।

অসমীয়া কিতাপৰ বাবে \texttt{language=assamese} বাছক; ডিফল্ট ভাষা বাংলা।
ফণ্টৰ কমাণ্ডবোৰ একেই থাকে; ৰ আৰু ৱ থকা ফণ্ট বাছি লওক। শিৰোনাম, নম্বৰ আৰু
ৰেফাৰেন্সৰ নাম বাছনি কৰা ভাষা অনুসৰি আহে।

লগত দিয়া \texttt{examples/bangla/main.tex}-ত এখন বাংলা কিতাপৰ শিৰোনাম পৃষ্ঠা,
সমল, কম্বিনেটৰিক্সৰ এটা সম্পূৰ্ণ অধ্যায়, অনুশীলনী আৰু তথ্যসূত্ৰ আছে। প্ৰথমে
LaTeX, তাৰ পিছত BibTeX, তাৰ পিছত আৰু দুবাৰ LaTeX চলাই কম্পাইল কৰক।

\section{শিৰোনাম পৃষ্ঠা আৰু অধ্যায়ৰ নক্সা}
শিৰোনাম পৃষ্ঠাটো memoir-ৰ চিনাকি \verb|\maketitle| কমাণ্ডেৰে তৈয়াৰ হয়;
কিতাপৰ নাম, লেখক আৰু তাৰিখ \verb|\title|, \verb|\author| আৰু \verb|\date|-ৰ
পৰা আহে। শিৰোনামৰ বাবে পৃথক পৃষ্ঠা ৰাখিবলৈ \texttt{titlingpage}
এনভায়ৰনমেণ্টৰ ভিতৰত \verb|\maketitle| লিখক। ক্লাছে নিজে শিৰোনাম পৃষ্ঠা
যোগ নকৰে। সম্পূৰ্ণ নিজৰ নক্সা বিচাৰিলে \verb|\maketitle| বাদ দি
\texttt{titlingpage}-ৰ ভিতৰত সেইটো লিখক। কপিৰাইটৰ লিখনিও আপোনাৰ
ডকুমেণ্টতেই ৰাখিব।

memoir-ৰ \verb|\chapterstyle| কমাণ্ডেৰে অধ্যায়ৰ নক্সা বাছি লওক।
ডিফল্ট নক্সা হৈছে \texttt{classic}। চাৰিওটা নক্সা ক'লা-বগা ছপাৰ উপযোগী;
memoir-ৰ আন নক্সাবোৰো ব্যৱহাৰ কৰিব পাৰে। memoir-ৰ সাধাৰণ অধ্যায়ৰ শিৰোনাম
পাবলৈ \verb|\chapterstyle{default}| লিখক। ইয়াৰ ফলত কেৱল অধ্যায়ৰ শিৰোনাম
সলনি হয়; পৃষ্ঠাৰ বিন্যাস বা পৰিচ্ছেদৰ নক্সা সলনি নহয়। সম্পূৰ্ণ পৃষ্ঠাত
নক্সাবোৰ চাবলৈ \texttt{examples/chapter-styles.pdf} খোলক। তলত একেটা অধ্যায়ৰ
শিৰোনাম চাৰিটা নক্সাত দেখুওৱা হৈছে; বক্সত ধৰাবলৈ আকাৰ সৰু কৰা হৈছে।
\begingroup
\let\chapter\manualchapterpreview
\begin{examplepair}
\chapterstyle{classic}
\chapter{গণনাৰ আৰম্ভণি}
\end{examplepair}
\begin{examplepair}
\chapterstyle{modern}
\chapter{গণনাৰ আৰম্ভণি}
\end{examplepair}
\begin{examplepair}
\chapterstyle{minimal}
\chapter{গণনাৰ আৰম্ভণি}
\end{examplepair}
\begin{examplepair}
\chapterstyle{ruled}
\chapter{গণনাৰ আৰম্ভণি}
\end{examplepair}
\endgroup

\subsection{নিজৰ মতে সজাওক}
তলৰ ছেটিংবোৰ সলনি কৰিবলৈ \verb|\bookbanglasetup{...}| ব্যৱহাৰ কৰক।
\begin{center}
\begin{tabular}{ll}
\toprule
ছেটিং & কাম\\
\midrule
\texttt{accent-color} & অধ্যায়ৰ নক্সাৰ ৰং (ডিফল্ট গাঢ় ধূসৰ)\\
\texttt{box-frame-color} & বক্সৰ সীমা আৰু শিৰোনাম (ডিফল্ট ক'লা)\\
\texttt{box-background-color} & বক্সৰ পটভূমি (ডিফল্ট বগা)\\
\texttt{quote-font} & উদ্ধৃতিৰ ফণ্টৰ কমাণ্ড (ডিফল্ট Chorus/হেলনীয়া)\\
\bottomrule
\end{tabular}
\end{center}
এইবোৰ BOOKbangla-ৰ নিজা ছেটিং। আন কিবা সলনি কৰিবলৈ ক্লাছ লোড কৰাৰ পিছত
memoir-ৰ কমাণ্ড ব্যৱহাৰ কৰক। পিছত দিয়া নিৰ্দেশ কাৰ্যকৰী হ'ব। উদাহৰণস্বৰূপে:
\begin{tcolorbox}[manualbox,title=\manualcodetitle]
\begin{code}
\chapterstyle{modern} % অধ্যায়ৰ নক্সা সলনি কৰক।
\pagestyle{plain}
\setlength{\parindent}{1em}
\setlength{\parskip}{0pt}
\end{code}
\end{tcolorbox}
ডিফল্ট নক্সালৈ উভতিবলৈ \verb|\chapterstyle{classic}| লিখক। পৰিচ্ছেদৰ
শিৰোনামৰ \S{} চিহ্নটো কেৱল সজ্জাৰ বাবে; ৰেফাৰেন্সত কেৱল পৰিচ্ছেদৰ নম্বৰ
থাকে। চিহ্নটো আঁতৰাবলৈ memoir-ৰ নম্বৰ সজোৱা কমাণ্ড ব্যৱহাৰ কৰক:
\begin{tcolorbox}[manualbox,title=\manualcodetitle]
\begin{code}
\setsecnumformat{\csname the#1\endcsname\quad}
\end{code}
\end{tcolorbox}
ডিফল্ট বিন্যাসত A4 কাগজ কাটি তৈয়াৰ কৰা পৃষ্ঠাৰ মাপ ৯.৫ বাই ৬.৮ ইঞ্চি।
memoir-এৰে কাটিবলৈ নিৰ্ধাৰণ কৰা মাপ, লিখনিৰ এলেকা আৰু মাৰ্জিন সলনি কৰি
\verb|\checkandfixthelayout| চলাওক। \verb|a5paper|-ৰ দৰে সৰু কাগজৰ অপশন
দিলে memoir-ৰ পৃষ্ঠাৰ মাপ আৰু মাৰ্জিন ব্যৱহাৰ হয়। কাগজ নকটাকৈ পোনে পোনে
A5 মাপৰ কিতাপৰ বাবে:
\begin{tcolorbox}[manualbox,title=\manualcodetitle]
\begin{code}
\setstocksize{210mm}{148mm}
\settrimmedsize{210mm}{148mm}{*}
\settrims{0pt}{0pt}
\setlrmarginsandblock{20mm}{18mm}{*}
\setulmarginsandblock{22mm}{25mm}{*}
\checkandfixthelayout
\end{code}
\end{tcolorbox}
যুগ্ম পৃষ্ঠাৰ ওপৰত কিতাপৰ নাম আৰু অযুগ্ম পৃষ্ঠাৰ ওপৰত পৰিচ্ছেদৰ নম্বৰ তথা
শিৰোনাম থাকে। memoir-ৰ পৃষ্ঠাশৈলীৰ কমাণ্ডেৰে এইবোৰ সলনি কৰিব পাৰে।
সবিশেষ \href{https://ctan.org/pkg/memoir}{memoir-ৰ নিৰ্দেশিকাত} পাব।

\section{উদাহৰণ, উপপাদ্য আৰু অনুশীলনী}
\texttt{example} হৈছে নম্বৰযুক্ত উদাহৰণৰ বক্স; প্ৰয়োজন হ'লে ই একাধিক
পৃষ্ঠালৈ বিয়পিব পাৰে। প্ৰতিটো অধ্যায়ত ইয়াৰ নম্বৰ পুনৰ আৰম্ভ হয়। সাধাৰণ
\verb|\label| কমাণ্ড আৰু \texttt{label} অপশন দুয়োটাই ব্যৱহাৰ কৰিব পাৰে:
\begin{examplepair}
\begin{example}[label=as:weekdays]
আঠজন বন্ধুৰ ভিতৰত অন্তত দুজনৰ জন্ম একেটা বাৰত কিয় হয়?
\end{example}
\begin{solution}
সপ্তাহত সাতটা বাৰ আছে।
\end{solution}
\autoref{as:weekdays} চোৱা।
\end{examplepair}
\texttt{solution}-এ প্ৰমাণৰ দৰে সমাধান সজায় আৰু শেষত সমাপ্তিচিহ্ন দিয়ে।
\verb|\begin{solution}[আন এটা সমাধান]| লিখি শিৰোনাম সলনি কৰিব পাৰে।
চিনাকি \texttt{proof} এনভায়ৰনমেণ্ট আৰু \verb|\qedhere| কমাণ্ডও ব্যৱহাৰ
কৰিব পাৰে। নম্বৰ নথকা মন্তব্য লিখিবলৈ \texttt{remark} ব্যৱহাৰ কৰক।

\texttt{reminder}-এ মনত ৰাখিবা বক্স তৈয়াৰ কৰে; ঐচ্ছিক আৰ্গুমেণ্টেৰে ইয়াৰ
শিৰোনাম সলনি কৰিব পাৰে। \texttt{diybox}-এ নিজে কৰা বক্স দিয়ে। উপপাদ্যৰ
বাবে \texttt{theorem} ব্যৱহাৰ কৰক। ইয়াত amsthm-ৰ সাধাৰণ \texttt{plain}
শৈলী অনুসৰি গাঢ় শিৰোনাম আৰু হেলনীয়া লিখনি থাকে; নম্বৰ অধ্যায়ভিত্তিক
পৃথক ক্ৰমত চলে। ঐচ্ছিক আৰ্গুমেণ্টত উপপাদ্যৰ নাম যোগ কৰিব পাৰে; চিনাকি লেবেল
আৰু ৰেফাৰেন্সৰ কমাণ্ডবোৰো কাম কৰে:
\begin{examplepair}
\begin{theorem}[কপৌৰ খোপ নীতি]
\label{as:theorem}
তিনিটা বস্তু দুটা খোপত ৰাখিলে
কোনো এটা খোপত অন্তত দুটা বস্তু থাকে।
\end{theorem}
\begin{proof}
প্ৰতিটো খোপত সৰ্বাধিক এটা বস্তু
ৰাখিলে দুটা বস্তু ৰাখিব পাৰি।
\end{proof}
\autoref{as:theorem} চোৱা।
\end{examplepair}
\begin{examplepair}
\begin{reminder}
বস্তু আৰু খোপ আগতে চিনাক্ত কৰা।
\end{reminder}
\begin{diybox}
চাৰিটা বস্তু লৈ চেষ্টা কৰা।
\end{diybox}
\end{examplepair}
অধিক উপপাদ্য বা অনুৰূপ এনভায়ৰনমেণ্টৰ প্ৰয়োজন হ'লে \verb|\newtheorem|-এৰে
তৈয়াৰ কৰক। সেইবোৰৰ নম্বৰ উপপাদ্যৰ সৈতে একে ক্ৰমত চলিব নে পৃথক হ'ব,
সেইটো আপুনি বাছি ল'ব পাৰে।

প্ৰতিটো \texttt{exercise}-ত এটা সমস্যা লিখিব; ইয়াৰ শিৰোনাম হয় সমস্যা।
সমস্যাবোৰ একেলগে ৰাখিবলৈ সাধাৰণ পৰিচ্ছেদ ব্যৱহাৰ কৰক:
\begin{tcolorbox}[manualbox,title=\manualcodetitle]
\begin{code}
\section{অনুশীলনী}
\end{code}
\end{tcolorbox}
\begin{examplepair}
\begin{exercise}[
  hint={মাহবোৰক খোপ বুলি ধৰা;
        $13>12$।}]
\label{as:months}
তেৰজন মানুহৰ ভিতৰত অন্তত দুজনৰ জন্ম একেটা মাহত হয়, দেখুওৱা।
\end{exercise}
\begin{exercise}
চাৰিটা পূৰ্ণসংখ্যাৰ ভিতৰত দুটাক
$3$-ৰে ভাগ কৰিলে ভাগশেষ সমান হয়, দেখুওৱা।
\end{exercise}
\autoref{as:months} চোৱা।
\end{examplepair}
ইংগিত দিলে সেইটো সমস্যাৰ পিছত ইংগিত: শিৰোনামেৰে দেখা যায়। নিদিলে কেৱল
সমস্যাটো ছপা হয়। সমস্যাৰ নম্বৰ উদাহৰণৰ নম্বৰৰ পৰা পৃথক; পৰিচ্ছেদ সলনি
হ'লেও ই চলি থাকে, নতুন অধ্যায়ত পুনৰ আৰম্ভ হয়। দুয়োটাতে \verb|\label|,
\verb|\ref| আৰু \verb|\autoref| ব্যৱহাৰ কৰিব পাৰে। \texttt{cleveref}-য়ো এই এনভায়ৰনমেণ্টবোৰৰ
অনূদিত নাম ব্যৱহাৰ কৰে। পৰিশিষ্টত memoir-ৰ অধ্যায়ৰ নম্বৰ ব্যৱহাৰ হয়, যিটো \LaTeX{}bangla-ই বাছনি কৰা ভাষাত দেখুৱায়।
প্ৰতিটো পৰিচ্ছেদত সমস্যাৰ নম্বৰ পুনৰ আৰম্ভ কৰিবলৈ প্ৰিঅ্যাম্বলত লিখক:
\begin{tcolorbox}[manualbox,title=\manualcodetitle]
\begin{code}
\counterwithin{exercise}{section}
\end{code}
\end{tcolorbox}
গোটেই কিতাপত সমস্যাৰ নম্বৰ একেৰাহে চলাবলৈ
\verb|\counterwithout{exercise}{chapter}| ব্যৱহাৰ কৰক। \texttt{example}
কাউণ্টাৰৰ ক্ষেত্ৰতো একেই কমাণ্ডবোৰ প্ৰযোজ্য।

\section{উদ্ধৃতি, তালিকা আৰু গণিত}
\begin{examplepair}
\begin{chapquote}[1em]{লেখক}
সৰু প্ৰশ্নৰ পৰাই ডাঙৰ ধাৰণাৰ জন্ম হয়।
\end{chapquote}
\end{examplepair}
ঐচ্ছিক আৰ্গুমেণ্টেৰে উদ্ধৃতিৰ দুয়োফালৰ খালী ঠাই নিৰ্ধাৰণ কৰিব পাৰে;
শেষত লেখকৰ নাম থাকে। ডিফল্ট হিচাপে ইংৰাজীত ১০৫\% স্কেলত TeX Gyre Chorus আৰু বাংলা বা
অসমীয়াত আপোনাৰ বাছনিৰ ফণ্টৰ হেলনীয়া ৰূপ ব্যৱহাৰ হয়। ইংৰাজীৰ ফণ্ট সলনি কৰিবলৈ:
\begin{tcolorbox}[manualbox,title=\manualcodetitle]
\begin{code}
\newfontfamily\quoteface{TeX Gyre Cursor}
\bookbanglasetup{quote-font={\quoteface\itshape}}
\end{code}
\end{tcolorbox}
উদ্ধৃতিৰ ভিতৰৰ বাংলা বা অসমীয়া লিখনি \LaTeX{}bangla-ই সামলায়। অনুচ্ছেদৰ আৰম্ভণিত
প্ৰথম আখৰটো ডাঙৰকৈ সজাব পাৰে:
\begin{examplepair}
\firstword{আ}{মি} চিনাকি এটা
ধাৰণাৰে আৰম্ভ কৰোঁ। সৰু প্ৰশ্নৰ
পৰা বহু ডাঙৰ ধাৰণাৰ জন্ম হয়।
\end{examplepair}

\texttt{smartlist}-ত ডিফল্টভাৱে হীৰাকৃতিৰ চিহ্ন থাকে। ইয়াত
\texttt{enumitem}-ৰ অপশনবোৰ ব্যৱহাৰ কৰি নিৰ্দিষ্ট তালিকাৰ চিহ্ন, মাৰ্জিন
আৰু ব্যৱধান সলনি কৰিব পাৰে:
\begin{examplepair}
\begin{smartlist}
\item বস্তুবোৰ চিনাক্ত কৰা।
\item খোপবোৰ গণনা কৰা।
\end{smartlist}
\begin{smartlist}[
  label=$\Rightarrow$,labelsep=.5em,
  leftmargin=2em,itemsep=0pt]
\item নিজৰ যুক্তি লিখি পেলোৱা।
\end{smartlist}
\end{examplepair}
চিনাকি \texttt{itemize} আৰু \texttt{enumerate} সাধাৰণভাৱেই কাম কৰে।

AMS-ৰ গণিত আৰু প্ৰমাণৰ সুবিধাবোৰ লোড কৰা আছে। বন্ধনীসহ সমীকৰণৰ ৰেফাৰেন্স
দিবলৈ \verb|\eqref| ব্যৱহাৰ কৰক। একাধিক শাৰীৰ সমীকৰণ পৃষ্ঠা পাৰ হৈ চলিব
পাৰে; সেইটো বন্ধ কৰিবলৈ \verb|\allowdisplaybreaks[0]| লিখক। ছবি যোগ কৰা
আৰু সাৰণীত পাদটীকা দিয়াও যায়। আন সুবিধাৰ বাবে memoir অথবা
\texttt{longtable}, \texttt{multicol}, \texttt{placeins}-ৰ দৰে পেকেজ ব্যৱহাৰ কৰক।

\section{তথ্যসূত্ৰ আৰু পেকেজৰ অপশন}
তথ্যসূত্ৰৰ পেকেজ, বিন্যাস আৰু তথ্যৰ ফাইল আপোনাৰ ডকুমেণ্টত বাছি লওক। ক্লাছে
\texttt{natbib} বা \texttt{biblatex} লোড নকৰে; উদাহৰণৰ কিতাপত
\texttt{biblatex} আৰু Biber ব্যৱহাৰ কৰা হৈছে। ক্লাছে \texttt{hyperref} লোড কৰে;
লিংকৰ ৰূপ সলনি কৰিবলৈ \verb|\hypersetup| ব্যৱহাৰ কৰক। \LaTeX{}bangla-ৰ
নম্বৰৰ অপশন সলনি কৰিবলৈ ক্লাছ লোড কৰাৰ আগতে সেইটো দিয়ক:
\begin{tcolorbox}[manualbox,title=\manualcodetitle]
\begin{code}
\PassOptionsToPackage{changecounternumbering=false}{latexbangla}
\documentclass[language=assamese]{bookbangla}
\end{code}
\end{tcolorbox}

\smallskip
কপিৰাইট ২০১৯, ২০২৬ আদীব হাসান। ৰক্ষণাবেক্ষণকাৰী আদীব হাসান।
ফাইলৰ তালিকা \texttt{README}-ত চাওক।
\end{document}
